% COLLECTION_METADATA: {"schema_version": 1, "kind": "reused_writeup", "independently_reviewed": false, "primary_source": "attacks/erdos/gpt_pro_5.2/809.tex", "source_paths": ["attacks/erdos/codex_5.2_extra_high/809.tex", "attacks/erdos/gpt_pro_5.2/809.tex"], "source_model": "gpt_pro_5.2", "source_completion": null, "source_claim": "unresolved"}
\section*{Erd\H{o}s Problem 809 --- collection record}

\subsection*{PROVENANCE AND STATUS}
Official problem: https://www.erdosproblems.com/809

Saved collaborative-database status: open.
Snapshot checked: 2026-09-23T17:22:24Z.
Latest status and formalization links: https://teorth.github.io/erdosproblems/

This is an attributed collection of existing writeups. The model-folder names below identify their original attribution; placement in GPT\_6\_Astra\_Ultra does not attribute their mathematical work to that model. All locally available source versions selected for this problem are reproduced below, without editing their text.

Proof claims, computations, literature statements, and completion estimates in the reused text have not been independently verified in this collection pass. A source claiming a solution does not change the saved upstream status.

Primary source (latest numbered version in the preferred model folder): \texttt{attacks/erdos/gpt\_pro\_5.2/809.tex}.

\subsection*{REUSED SOURCE 1}
Original attribution: \texttt{codex\_5.2\_extra\_high}.
Original file: \texttt{attacks/erdos/codex\_5.2\_extra\_high/809.tex}.

% BEGIN REUSED SOURCE: attacks/erdos/codex_5.2_extra_high/809.tex
\noindent\textbf{1) FORMAL RESTATEMENT.}

Fix an integer $k\ge 3$. For each integer $n\ge 1$, let
\[
 m(n):=\left\lfloor \frac{n^2}{4}\right\rfloor+1.
\]
Define $F_k(n)$ to be the minimal integer $r\ge 1$ such that there exist
\begin{itemize}
\item a simple graph $G$ on $n$ vertices with exactly $m(n)$ edges, and
\item an edge-coloring $\chi: E(G)\to [r]:=\{1,\dots,r\}$
\end{itemize}
with the property:

\emph{(Rainbow $C_{2k+1}$ property)} For every subgraph $H\subseteq G$ isomorphic to the cycle $C_{2k+1}$, the $2k+1$ edges of $H$ receive pairwise distinct colors under $\chi$.

Question: Is it true that
\[
F_k(n)\sim \frac{n^2}{8}\qquad (n\to\infty)?
\]

\medskip
\noindent\textbf{2) QUICK LITERATURE/CONTEXT CHECK.}

The problem text states that Burr--Erd\H{o}s--Graham--S\'os proved $F_k(n)\gg n^2$ (no proof given here). I do not use any additional literature.

\medskip
\noindent\textbf{3) ATTACK PLAN.}

\emph{Proof (upper bound) strategies.}
\begin{itemize}
\item Construct a specific $n$-vertex graph with $m(n)$ edges and an explicit coloring using about $n^2/8$ colors such that every $C_{2k+1}$ is rainbow.
\item Start from an extremal bipartite graph $K_{\lfloor n/2\rfloor,\lceil n/2\rceil}$ (which has $\lfloor n^2/4\rfloor$ edges) and add one edge; understand the resulting family of $C_{2k+1}$'s and design a near-optimal coloring for those cycles.
\end{itemize}

\emph{Disproof/lower bound strategies.}
\begin{itemize}
\item Prove that in any such $G$ with $m(n)$ edges, any coloring with $o(n^2)$ colors forces a repeated color on some $C_{2k+1}$ by a counting/double-counting argument over cycles.
\item Seek a graph $G$ with $m(n)$ edges and relatively few (or specially structured) $C_{2k+1}$'s, potentially allowing a coloring with $o(n^2)$ colors, which would refute the conjectured constant $1/8$.
\end{itemize}

I did not obtain a proof or counterexample; below are fully proved elementary bounds and structural observations.

\medskip
\noindent\textbf{4) WORK.}

\textbf{Lemma 809.1 (trivial universal upper bound).}
For every $k\ge 3$ and $n\ge 1$,
\[
F_k(n)\le m(n)=\left\lfloor \frac{n^2}{4}\right\rfloor+1.
\]

\emph{Proof.}
Take any graph $G$ on $n$ vertices with $m(n)$ edges (such graphs exist for all $n\ge 3$). Color every edge with a distinct color. Then every subgraph, in particular every $C_{2k+1}$, has all edges different colors. This uses exactly $m(n)$ colors, so $F_k(n)\le m(n)$.
\qed

\textbf{Lemma 809.2 (odd cycles in a near-bipartite extremal graph).}
Let $a,b\ge 1$ and let $G$ be obtained from the complete bipartite graph $K_{a,b}$ (with bipartition $A\cup B$) by adding one extra edge $uv$ inside $A$.
Then for every integer $j$ with
\[
1\le j\le \min\{a-1,b\},
\]
$G$ contains a cycle of length $2j+1$.

\emph{Proof.}
Fix such a $j$. Choose distinct vertices
\[
 u=v_0,\ v=v_{2j}\in A,
\]
and choose further distinct vertices $v_2,v_4,\dots,v_{2j-2}\in A\setminus\{u,v\}$ (this is possible because $A$ has size $a$ and we need $j-1\le a-2$ additional vertices, i.e. $j\le a-1$).
Also choose distinct vertices $v_1,v_3,\dots,v_{2j-1}\in B$ (possible because we need $j\le b$ vertices in $B$).

Now consider the cyclic sequence
\[
 v_0(=u),\ v_1,\ v_2,\ v_3,\dots,\ v_{2j-1},\ v_{2j}(=v),\ v_0.
\]
For each even index $2i$, $v_{2i}\in A$ and $v_{2i+1}\in B$, so the edge $v_{2i}v_{2i+1}$ lies in $K_{a,b}$. Similarly $v_{2i+1}v_{2i+2}$ lies in $K_{a,b}$.
Finally, the closing edge $v_{2j}v_0$ is exactly the added edge $uv$ inside $A$.
All vertices in the sequence are distinct by construction, so these edges form a simple cycle of length $2j+1$.
\qed

\textbf{FAST REALITY CHECK (tiny instances by computation).}

Two basic sanity checks:
\begin{itemize}
\item If $n<2k+1$ then no graph on $n$ vertices contains a $C_{2k+1}$, so the rainbow-cycle condition is vacuous and $F_k(n)=1$.
\item Even when $n=2k+1$, it can happen that $F_k(n)=1$ because one can choose $G$ with $m(n)$ edges but no $C_{2k+1}$.
\end{itemize}

Concrete example for $k=3$ and $n=7$: here $m(7)=\lfloor 49/4\rfloor+1=13$.
Let $G$ be the disjoint union of an isolated vertex and a $6$-vertex graph formed from $K_6$ by deleting two edges. Then $|E(G)|=13$ and $G$ has no $C_7$ (because every component has at most $6$ vertices).
A direct programmatic check confirms ``has $C_7$? False'' for this $G$. Therefore $F_3(7)=1$.

\medskip
\noindent\textbf{5) VERIFICATION.}

\begin{itemize}
\item Lemma 809.1 is immediate; unique-edge coloring always enforces rainbow-ness on any fixed cycle.
\item Lemma 809.2: verified that the constructed cycle uses only edges from $K_{a,b}$ except for the single added intra-$A$ edge, and that all vertices are distinct.
\item The $k=3,n=7$ sanity-check graph is disconnected with maximum component size $6$, so it cannot contain a $7$-cycle; the computation confirms this.
\end{itemize}

\medskip
\noindent\textbf{6) FINAL.} \textbf{UNRESOLVED.}

\begin{enumerate}
\item[(i)] \emph{Strongest proved partial result here:} the trivial upper bound $F_k(n)\le \lfloor n^2/4\rfloor+1$ (Lemma 809.1) and a structural fact that the natural extremal ``$K_{\lfloor n/2\rfloor,\lceil n/2\rceil}$ plus one edge'' construction contains many odd cycles of all lengths $3,5,\dots,2\min\{\lfloor n/2\rfloor-1,\lceil n/2\rceil\}+1$ (Lemma 809.2).
\item[(ii)] \emph{First gap (crisp):} either construct, for each large $n$, a graph with $m(n)$ edges admitting a rainbow-$C_{2k+1}$ edge-coloring using $(1/8+o(1))n^2$ colors, or prove that any such coloring requires at least $(1/8-o(1))n^2$ colors.
\item[(iii)] \emph{Top 3 next moves:}
  \begin{enumerate}
  \item For the candidate graph $K_{\lfloor n/2\rfloor,\lceil n/2\rceil}$ plus one intra-part edge, explicitly characterize all $C_{2k+1}$ cycles and attempt to color the bipartite edges so that no two edges of the same color can lie on any such cycle.
  \item Prove a lower bound by counting $C_{2k+1}$'s and showing that if a color class contains too many edges then two of them must co-occur on some $C_{2k+1}$.
  \item Run exact search for the smallest nontrivial instances (e.g. $n=2k+2$ or $2k+3$) using SAT/ILP to determine $F_k(n)$ and to guess extremal structures.
  \end{enumerate}
\item[(iv)] \emph{What a minimal counterexample would likely look like:} if the asymptotic constant $1/8$ is false, then there should exist graphs with $m(n)$ edges for which the family of $C_{2k+1}$'s is structured so that edges can be grouped into color classes of average size substantially larger than $2$ without creating two same-colored edges on a $C_{2k+1}$; equivalently, a graph where the ``co-occurrence graph'' on edges (joining two edges if they lie on a common $C_{2k+1}$) has unusually small chromatic number.
\end{enumerate}



% END REUSED SOURCE: attacks/erdos/codex_5.2_extra_high/809.tex

\subsection*{REUSED SOURCE 2}
Original attribution: \texttt{gpt\_pro\_5.2}.
Original file: \texttt{attacks/erdos/gpt\_pro\_5.2/809.tex}.

% BEGIN REUSED SOURCE: attacks/erdos/gpt_pro_5.2/809.tex
## 1) FORMAL RESTATEMENT

Fix an integer (k\ge 3). For each integer (n\ge 1), define (F_k(n)) to be the minimum integer (r\ge 1) such that there exist

* a **simple** undirected graph (G=(V,E)) with (|V|=n) and (|E|=\lfloor n^2/4\rfloor+1), and
* an **edge-colouring** (\varphi:E\to [r]={1,2,\dots,r}),

with the property that **every** subgraph (H\subseteq G) isomorphic to the odd cycle (C_{2k+1}) is **rainbow** under (\varphi); i.e. (\varphi) is injective on (E(H)) (no colour repeats among the (2k+1) edges of that cycle).

Question (Burr–Erdős–Graham–Sós):
[
F_k(n)\sim \frac{n^2}{8}\quad (n\to\infty),
]
meaning (\displaystyle \lim_{n\to\infty}\frac{F_k(n)}{n^2}=\frac18).

Conventions/edge cases:

* If (n<2k+1), then (G) contains no (C_{2k+1}), so the rainbow constraint is vacuous.
* Logs are irrelevant here.
* Graphs are finite, simple, undirected.

Stress points:

* The edge count (\lfloor n^2/4\rfloor+1) is just **one above** Mantel’s bound for triangle-free / bipartite extremality, so (G) is “barely non-bipartite”.
* A single (C_{2k+1}) already forces (r\ge 2k+1).

---

## 2) QUICK LITERATURE/CONTEXT CHECK (web browsing)

* The ErdosProblems entry for #809 states: Burr–Erdős–Graham–Sós proved (F_k(n)\gg n^2).
* For context: for (n) large compared to (k), the extremal number for odd cycles satisfies
  [
  \mathrm{ex}(n,C_{2k+1})=\left\lfloor\frac{n^2}{4}\right\rfloor \quad\text{(for } n\ge 4k-2\text{)},
  ]
  attributed to Füredi–Gunderson.
  Hence for (n\ge 4k-2), any graph with (\lfloor n^2/4\rfloor+1) edges must contain at least one (C_{2k+1}), so the rainbow condition is asymptotically non-vacuous.

I did not find any indication (up to the site’s 2025 update cadence) that the asymptotic (\sim n^2/8) has been resolved.

---

## 3) ATTACK PLAN

**Proof strategies (to try to prove (F_k(n)\sim n^2/8))**

1. **Upper bound construction:** build (G) with (\lfloor n^2/4\rfloor+1) edges and an edge-colouring using ((1/8+o(1))n^2) colours. Natural route: make the “cycle-conflict graph” on edges sparse enough that many edges can share colours safely (no two same-coloured edges lie on a common (C_{2k+1})).
2. **Lower bound via conflicts:** show that in every such (G), there is a large set of edges pairwise forced to have distinct colours because any two lie together in some (C_{2k+1}). Turning this into a sharp constant (1/8) is the hard part.
3. **Stability + supersaturation:** use that (G) is near-bipartite (by stability around Mantel) and argue that the forced “odd-cycle mechanism” creates a dense conflict relation on at least (\sim n^2/4) edges, yielding a quadratic colour lower bound with a constant.

**Disproof/construction strategies (to falsify (\sim n^2/8))**

1. Produce a family with (F_k(n)\le (1/8-\varepsilon)n^2) (a better colouring than conjectured).
2. Prove a universal lower bound (F_k(n)\ge (1/8+\varepsilon)n^2) (show conjectured constant too small).

**Best path chosen in the work below:** I focused on (i) small-case sanity and (ii) analyzing the most natural “one edge above bipartite extremal” graph, to see what colour lower bounds it forces. This yields a rigorous (but negative-for-construction) structural lemma.

---

## 4) WORK

### PHASE 1 — FAST REALITY CHECK

**Vacuous regime.** If (n<2k+1), then no (C_{2k+1}) exists in any (n)-vertex graph, so the rainbow condition is vacuous. Any graph (G) with (\lfloor n^2/4\rfloor+1) edges (exists for all (n\ge 3)) can be coloured with one colour. Hence:
[
\boxed{n<2k+1\ \Rightarrow\ F_k(n)=1.}
]

**Non-vacuous regime for large (n).** If (n\ge 4k-2), then (\mathrm{ex}(n,C_{2k+1})=\lfloor n^2/4\rfloor).
So any (G) with (\lfloor n^2/4\rfloor+1) edges contains at least one (C_{2k+1}), and a rainbow (C_{2k+1}) needs at least (2k+1) distinct colours. Therefore:
[
\boxed{n\ge 4k-2\ \Rightarrow\ F_k(n)\ge 2k+1.}
]

Trivial upper bound:
[
F_k(n)\le \lfloor n^2/4\rfloor+1
]
by colouring each edge distinctly.

---

### PHASE 2 — A rigorous “stress test” on the obvious near-extremal graph

A natural candidate (G) with (\lfloor n^2/4\rfloor+1) edges is:
[
G_0 := K_{a,b}\ \cup\ {uv},
]
where (a=\lfloor n/2\rfloor), (b=\lceil n/2\rceil), (A\sqcup B) is the bipartition of (K_{a,b}), and (u,v\in A) are distinct.

This is “one extra edge on top of Mantel”.

I now prove that in (G_0), the rainbow requirement actually forces **almost all** bipartite edges to have **distinct** colours, i.e. (r) must be (\sim n^2/4). This shows (G_0) cannot be close to achieving an (n^2/8) upper bound.

#### Lemma (two bipartite edges are co-cyclic with (uv))

Fix (k\ge 3). Suppose (a\ge k+1) and (b\ge k). Let (G_0=K_{a,b}\cup{uv}) as above with (u,v\in A).

Let (e_1,e_2) be two distinct edges of the complete bipartite part with endpoints in ((A\setminus{u,v})\times B). Then there exists a cycle (C\subseteq G_0) with
[
C\cong C_{2k+1}\quad\text{and}\quad {uv,e_1,e_2}\subseteq E(C).
]

**Proof.** It suffices to construct a **simple** (u)-to-(v) path (P) of length (2k) inside (K_{a,b}) that contains both (e_1) and (e_2). Then adding the edge (uv) closes (P) into a (C_{2k+1}).

Because the bipartite part is complete, any alternating sequence of distinct vertices
[
u\in A,\ B,\ A,\ B,\ \dots,\ B,\ v\in A
]
gives a valid path. So we only must ensure:

1. alternation,
2. no vertex repetition,
3. the two required edges appear somewhere among the consecutive pairs.

Since (k\ge 3), the path has at least (k) (B)-vertices and (k-1) intermediate (A)-vertices. Our assumptions (b\ge k), (a\ge k+1) guarantee there are enough distinct vertices available to realize any prescribed constant-size “pattern” near the start (containing (e_1,e_2)) and then complete the remainder with fresh vertices.

Concretely: build an initial alternating prefix starting at (u) that traverses (e_1) and (e_2) (if they share an endpoint, place them consecutively through that shared endpoint; if they are disjoint, place them as two separate consecutive (B!-!A) steps). This uses at most 2 intermediate (A)-vertices from (A\setminus{u,v}) and at most 2 vertices from (B). Then extend the alternating walk to reach (v) using fresh vertices; completeness ensures all required edges exist. The result is a simple (u)-to-(v) path of length exactly (2k) containing (e_1,e_2). Adding (uv) yields the desired (C_{2k+1}). ∎

#### Corollary (large forced number of colours in (G_0))

In any edge-colouring (\varphi) of (G_0) in which every (C_{2k+1}) is rainbow, all edges between (A\setminus{u,v}) and (B) must have **pairwise distinct** colours. Hence
[
r \ \ge\ |E(A\setminus{u,v},B)|=(a-2)b.
]
For balanced (a=\lfloor n/2\rfloor), (b=\lceil n/2\rceil),
[
(a-2)b=\left\lfloor\frac{n^2}{4}\right\rfloor-\Theta(n),
]
so this particular (G_0) requires
[
\boxed{r\ge \left\lfloor\frac{n^2}{4}\right\rfloor-O(n).}
]

**Proof.** If two distinct edges (e_1,e_2) in (E(A\setminus{u,v},B)) shared a colour, then by the lemma there is a (C_{2k+1}) containing both, so that cycle would repeat a colour—contradiction. Therefore all those edges have distinct colours, so (r) is at least their count. ∎

**Interpretation:** the “obvious” near-extremal graph (K_{a,b}+uv) is extremely hostile to colour reuse; it forces (\sim n^2/4) colours, not (\sim n^2/8). So any graph witnessing (F_k(n)\sim n^2/8) (if true) must have a much sparser “two edges lie on a common (C_{2k+1})” relation on a positive fraction of its edges.

---

## 5) VERIFICATION (adversarial self-check)

* **Vacuous case (n<2k+1):** correct; no (C_{2k+1}) subgraph exists.
* **Non-vacuous trigger (n\ge 4k-2):** depends on (\mathrm{ex}(n,C_{2k+1})=\lfloor n^2/4\rfloor) for that range, cited.
* **Key lemma on (G_0):** the only delicate point is ensuring the two specified edges can be embedded in a simple alternating (u)-to-(v) path of length (2k) without vertex repetition. The availability conditions (a\ge k+1), (b\ge k) ensure we can choose all intermediate vertices distinct while reserving the endpoints needed for the forced edges. Completeness of (K_{a,b}) guarantees the necessary adjacencies.

No step uses hidden assumptions beyond simplicity and completeness.

---

## 6) FINAL: **UNRESOLVED**

I did **not** produce either:

* a complete proof that (F_k(n)\sim n^2/8), or
* an explicit counterexample disproving that asymptotic.

### (i) Strongest fully proved partial results here

1. (F_k(n)=1) for (n<2k+1).
2. For (n\ge 4k-2), (F_k(n)\ge 2k+1).
3. For the canonical graph (G_0=K_{\lfloor n/2\rfloor,\lceil n/2\rceil}+uv), any valid colouring needs
   [
   r\ge (a-2)b = \left\lfloor\frac{n^2}{4}\right\rfloor - O(n),
   ]
   so (G_0) cannot witness a (\sim n^2/8) upper bound.

### (ii) First gap

Construct (or rule out) an (n)-vertex graph with (\lfloor n^2/4\rfloor+1) edges whose edge-set can be partitioned into (\sim n^2/8) colour classes such that no two edges in a class lie on a common (C_{2k+1}).

### (iii) Top 3 next moves

1. Identify minimizers of (F_k(n)) structurally: what graphs with (\lfloor n^2/4\rfloor+1) edges minimize the “cycle conflict graph” on edges?
2. Attempt an explicit upper-bound construction where every (C_{2k+1}) must pass through a sparse “bottleneck” set of edges (so colours can be reused heavily outside that set).
3. Compute small-(n) optimizers (fixed (k)) to guess the extremal pattern for (G) and for the colouring (even a heuristic can suggest the right structure).

### (iv) What a minimal counterexample would likely look like

If the conjecture (F_k(n)\sim n^2/8) is false, then for some (\varepsilon>0) infinitely many (n) would satisfy either:

* (F_k(n)\le (1/8-\varepsilon)n^2) via a graph where most edges fall into “non-co-cyclic” pairs (allowing more reuse), or
* (F_k(n)\ge (1/8+\varepsilon)n^2) via a universal argument that every such (G) has a conflict relation too dense to permit (\sim m/2) colouring.

My proved lemma shows the most obvious near-Mantel graph (K_{a,b}+uv) lies in the *opposite* extreme, forcing (\sim n^2/4) colours.

% END REUSED SOURCE: attacks/erdos/gpt_pro_5.2/809.tex

\subsection*{COMPLETION ESTIMATE}
The primary source supplies no extractable numerical completion estimate. Its mathematical progress is not numerically assessed here.

Newly verified progress in this collection pass: $0\%$. This measures new verification work, not the amount of mathematics achieved in the reused text.
